\chapter{Annales Réelles et Décorticage des Sujets de Dissertation}

Ce chapitre présente l'analyse approfondie des sujets \textbf{réellement posés} lors des récentes sessions du concours d'Ingénieurs d'État du Ministère de l'Économie et des Finances, ainsi que les sujets à très haute probabilité issus des réformes majeures de l'État.

Chaque sujet est traité selon la matrice exigée :
\begin{center}
\textbf{Sujet $\longrightarrow$ Problématique $\longrightarrow$ Plan détaillé $\longrightarrow$ Idées principales $\longrightarrow$ Exemples marocains $\longrightarrow$ Conclusion}
\end{center}

% =========================================================================
\section{Sujet 1 (Session Réelle du 28 Avril 2024) : Digitalisation et Stratégie « Maroc Digital 2030 »}
% =========================================================================

\begin{tcolorbox}[colback=blue!5!white,colframe=blue!75!black,title=\textbf{Énoncé Officiel du Sujet (28 Avril 2024)}]
\textit{« Durant les dernières années, le Maroc a fait beaucoup de réalisations dans le domaine de la digitalisation. Quel bilan faites-vous de ces réalisations ? Quelles sont les axes prioritaires de la Stratégie Maroc Digital 2030 et les orientations stratégiques pour accompagner cette transition ? »}
\end{tcolorbox}

\subsection*{1. Problématique}
Dans quelle mesure les acquis de la transition numérique marocaine constituent-ils un levier de modernisation économique et administrative, et comment la stratégie « Maroc Digital 2030 » entend-elle surmonter les fractures structurelles pour faire du digital un vecteur de souveraineté et d'inclusion ?

\subsection*{2. Plan Détaillé Proposé}
\begin{itemize}
    \item \textbf{I. Le Bilan de la Digitalisation au Maroc : Des Avancées Notables Face à des Freins Persistants}
    \begin{itemize}
        \item \textbf{A.} Des réalisations substantielles dans les services publics et l'offshoring.
        \item \textbf{B.} Des limites structurelles : fracture territoriale, compétences et inclusion des TPME.
    \end{itemize}
    \item \textbf{II. La Stratégie « Maroc Digital 2030 » et les Leviers d'Accompagnement de la Transition}
    \begin{itemize}
        \item \textbf{A.} Les deux piliers et catalyseurs prioritaires de la feuille de route 2030.
        \item \textbf{B.} Les orientations stratégiques d'accompagnement : souveraineté des données, 5G, talents et gouvernance.
    \end{itemize}
\end{itemize}

\subsection*{3. Idées Principales et Arguments}
\begin{itemize}
    \item \textbf{Pour le bilan (Partie I) :}
    \begin{itemize}
        \item Émergence du Maroc comme 2\textsuperscript{e} hub d'outsourcing en Afrique (plus de 130 000 emplois créés, exportations de services IT en forte croissance).
        \item Révolution des services publics dématérialisés : téléprocédures fiscales (SIMPL de la DGI), douanières (BADR de l'ADII), marchés publics électroniques (TGR), et simplification administrative (Loi 55-19 et portail national \textit{Idarati}).
        \item Limites : persistance de la fracture numérique entre zones urbaines et rurales, taux d'adoption digital encore faible chez les TPME traditionnelles, déficit en compétences pointues (fuite des cerveaux IT), et manque d'interopérabilité fluide entre ministères.
    \end{itemize}
    \item \textbf{Pour Maroc Digital 2030 (Partie II) :}
    \begin{itemize}
        \item \textbf{Pilier 1 : Digitalisation des services publics} (atteindre un parcours 100\% digitalisé et ergonomique, guichet unique, interopérabilité généralisée).
        \item \textbf{Pilier 2 : Dynamisation de l'économie numérique} (développer les startups, exporter des solutions numériques locales, attirer les investissements étrangers dans la tech).
        \item \textbf{Catalyseurs clés :} Déploiement de la 5G, développement d'un Cloud souverain national, formation de 100 000 jeunes talents par an d'ici 2030, intégration de l'Intelligence Artificielle dans l'analyse de données fiscales et financières.
    \end{itemize}
\end{itemize}

\subsection*{4. Exemples et Données Marocaines Précises}
\begin{itemize}
    \item Budget alloué à la première phase (2024-2026) : environ 11 milliards de dirhams.
    \item Cadre juridique : Loi n° 09-08 (protection des données personnelles via la CNDP), Loi n° 43-20 (services de confiance pour les transactions électroniques), et Loi n° 55-19 (simplification des procédures administratives).
    \item Direction Générale de la Sécurité des Systèmes d'Information (DGSSI) et Directive Nationale de la Sécurité des SI (DNSSI).
    \item Écosystème startups : Fonds InnovX, programmes de la Caisse des Dépôts et de Gestion (CDG Invest).
\end{itemize}

\subsection*{5. Conclusion et Ouverture}
Le Maroc passe d'une numérisation d'appoint à une véritable économie de la donnée. L'enjeu futur réside dans la souveraineté numérique (Data Centers locaux) et l'utilisation de l'intelligence artificielle pour optimiser les recettes fiscales et lutter contre l'économie informelle.

% =========================================================================
\section{Sujet 2 (Session Réelle du 28 Avril 2024) : Transition et Sécurité Énergétique au Maroc}
% =========================================================================

\begin{tcolorbox}[colback=blue!5!white,colframe=blue!75!black,title=\textbf{Énoncé Officiel du Sujet (28 Avril 2024)}]
\textit{« À l'aune de la crise énergétique que traverse l'économie mondiale, l'accélération de la transition énergétique est placée au centre des priorités du développement économique de notre pays. Quels sont les nouveaux enjeux de la sécurité énergétique au Maroc ? Quels sont les axes stratégiques pour renforcer cette sécurité et les défis liés à la stratégie nationale de transition énergétique ? »}
\end{tcolorbox}

\subsection*{1. Problématique}
Face aux chocs géopolitiques et à la volatilité des marchés énergétiques internationaux, comment le Maroc peut-il concilier l'impératif immédiat de sécurisation de ses approvisionnements avec l'ambition durable d'une transition énergétique basée sur les énergies renouvelables et l'hydrogène vert ?

\subsection*{2. Plan Détaillé Proposé}
\begin{itemize}
    \item \textbf{I. Les Nouveaux Enjeux de la Sécurité Énergétique Marocaine : Vulnérabilités et Défis}
    \begin{itemize}
        \item \textbf{A.} La dépendance énergétique externe et l'impact direct sur les finances publiques.
        \item \textbf{B.} L'impératif de décarbonation industrielle et les nouvelles exigences climatiques mondiales.
    \end{itemize}
    \item \textbf{II. Les Axes Stratégiques et les Leviers Opérationnels de la Transition Énergétique}
    \begin{itemize}
        \item \textbf{A.} L'accélération du mix renouvelable et le déploiement de « l'Offre Maroc » pour l'Hydrogène Vert.
        \item \textbf{B.} Efficacité énergétique, interconnexions régionales et défis de financement.
    \end{itemize}
\end{itemize}

\subsection*{3. Idées Principales et Données Marocaines}
\begin{itemize}
    \item Taux de dépendance énergétique du Maroc : passé de plus de 98\% au début des années 2000 à environ 90\% aujourd'hui, avec un coût lourd sur la balance commerciale et les charges de compensation.
    \item Objectif national fixé par les Orientations Royales : dépasser 52\% de capacité électrique installée d'origine renouvelable (solaire, éolien, hydraulique) avant 2030.
    \item Réalisations majeures : Complexe solaire Noor Ouarzazate (580 MW), parcs éoliens de Tarfaya et Midelt, stations de transfert d'énergie par pompage (STEP Abdelmoumen).
    \item « Offre Maroc » pour l'hydrogène vert (circulaire du Chef du Gouvernement de mars 2024) : mobilisation de 1 million d'hectares de foncier public pour attirer les investisseurs mondiaux.
    \item Rôle du MEF : financement vert, partenariats public-privé (loi 86-12), et décarbonation pour anticiper le Mécanisme d'Ajustement Carbone aux Frontières (MACF) de l'Union Européenne.
\end{itemize}

\subsection*{4. Conclusion}
La transition énergétique marocaine n'est plus seulement un choix écologique, mais une exigence de compétitivité industrielle et de souveraineté budgétaire.

% =========================================================================
\section{Sujet 3 (Sessions Réelles 2022-2023) : L'Inflation et la Préservation du Pouvoir d'Achat}
% =========================================================================

\begin{tcolorbox}[colback=blue!5!white,colframe=blue!75!black,title=\textbf{Énoncé Réel / Thème Fréquent}]
\textit{« L'inflation au Maroc s'est nettement accélérée ces dernières années. Quelles sont les principales causes de cette hausse continue des prix ? Quelles sont ses répercussions économiques et sociales et quelles sont les mesures de politique économique et budgétaire adoptées ou recommandées pour y faire face ? »}
\end{tcolorbox}

\subsection*{1. Problématique}
Quelles sont les dynamiques internes et externes expliquant la flambée inflationniste récente au Maroc, et quelle articulation entre politique monétaire (Bank Al-Maghrib) et politique budgétaire (MEF) permet de contenir la hausse des prix sans étouffer la croissance ?

\subsection*{2. Plan Détaillé}
\begin{itemize}
    \item \textbf{I. Diagnostic de l'Inflation : Causes Multiples et Impacts Socio-Économiques}
    \begin{itemize}
        \item \textbf{A.} Les facteurs déclencheurs : inflation importée (énergie, matières premières) et facteurs climatiques (sécheresse sévère affectant les prix agricoles).
        \item \textbf{B.} Les conséquences : érosion du pouvoir d'achat des ménages, hausse des coûts de production et tension sur le budget de l'État.
    \end{itemize}
    \item \textbf{II. Les Réponses des Politiques Publiques et les Perspectives de Stabilisation}
    \begin{itemize}
        \item \textbf{A.} La politique monétaire de Bank Al-Maghrib : relèvements successifs du taux directeur.
        \item \textbf{B.} Les mesures budgétaires du MEF et réformes structurelles : soutien aux intrants agricoles, subventions ciblées et réorganisation des circuits de commercialisation.
    \end{itemize}
\end{itemize}

\subsection*{3. Chiffres et Faits Marocains Clés}
\begin{itemize}
    \item Pic d'inflation en 2022-2023 (6,6\% en 2022, 6,1\% en 2023), tirée principalement par les produits alimentaires (+15\% à son pic), avant une décélération progressive en 2024 vers 1,5\% à 2\%.
    \item Décisions de Bank Al-Maghrib : augmentation graduelle du taux directeur de 1,50\% jusqu'à 3,00\% pour freiner les anticipations inflationnistes, puis premier assouplissement à 2,75\% en juin 2024 face à l'accalmie des prix.
    \item Mesures MEF : soutien direct aux professionnels du transport routier pour éviter la répercussion sur les prix à la consommation, suspension des droits de douane et de la TVA sur l'importation de bétail et d'huile, et enveloppes exceptionnelles de la Caisse de Compensation (gaz butane, blé, sucre).
\end{itemize}

% =========================================================================
\section{Sujet 4 (Session Réelle 2021-2022) : Maîtrise du Déficit Budgétaire et Dette Publique}
% =========================================================================

\begin{tcolorbox}[colback=blue!5!white,colframe=blue!75!black,title=\textbf{Énoncé Réel}]
\textit{« Le déficit des finances publiques constitue une caractéristique structurelle de l'économie marocaine. Quelles sont, à votre avis, les mesures à adopter pour maîtriser l'impact de ce déficit tout en préservant la dynamique d'investissement et les impératifs sociaux ? »}
\end{tcolorbox}

\subsection*{1. Problématique}
Comment l'État marocain peut-il réussir l'équation complexe de la consolidation budgétaire (réduction du déficit sous 4\% du PIB) tout en finançant les chantiers stratégiques de l'État social et les grands investissements d'infrastructures ?

\subsection*{2. Plan Détaillé}
\begin{itemize}
    \item \textbf{I. La Problématique du Déficit Budgétaire et de la Dette : Facteurs et Contraintes}
    \begin{itemize}
        \item \textbf{A.} Les causes structurelles du déficit : rigidité des dépenses obligatoires (masse salariale, charge de la dette, compensation).
        \item \textbf{B.} Le risque d'éviction et les seuils d'alerte de l'endettement du Trésor ($\approx 70\%$ du PIB).
    \end{itemize}
    \item \textbf{II. La Stratégie d'Assainissement Budgétaire et les Nouveaux Leviers de Financement}
    \begin{itemize}
        \item \textbf{A.} L'élargissement de l'assiette fiscale et la rationalisation des dépenses (LOF 130-13 et Loi-cadre 69-19).
        \item \textbf{B.} Les mécanismes innovants de financement et le partenariat public-privé (Fonds Mohammed VI pour l'Investissement, ANGSPE).
    \end{itemize}
\end{itemize}

\subsection*{3. Chiffres et Références Institutionnelles}
\begin{itemize}
    \item Évolution du déficit budgétaire : de plus de 7\% du PIB en 2020 (crise Covid) à 4,3\% en 2023, avec pour cible 4\% en 2024 et 3,5\% à l'horizon 2025-2026.
    \item Dette du Trésor : stabilisée autour de 69,5\% à 70\% du PIB, avec une composante intérieure prédominante (environ 75\% de la dette) limitant le risque de change.
    \item Financements innovants : monétisation d'actifs immobiliers de l'État par la Direction des Domaines et gestion active du portefeuille d'entreprises publiques.
\end{itemize}

% =========================================================================
\section{Sujet 5 (Thème Royal Majeur) : Généralisation de la Protection Sociale}
% =========================================================================

\begin{tcolorbox}[colback=blue!5!white,colframe=blue!75!black,title=\textbf{Thème Stratégique et Sujet Très Probable}]
\textit{« Le chantier de la généralisation de la protection sociale au Maroc : réalisations, modalités de financement, soutenabilité budgétaire et défis de mise en œuvre. »}
\end{tcolorbox}

\subsection*{1. Problématique}
En quoi la généralisation de la protection sociale représente-t-elle un changement de paradigme vers un véritable État social, et quelles sont les conditions de sa pérennité budgétaire et opérationnelle ?

\subsection*{2. Plan Détaillé}
\begin{itemize}
    \item \textbf{I. Un Projet Sociétal d'Envergure : Piliers et Réalisations Concrètes}
    \begin{itemize}
        \item \textbf{A.} La feuille de route de la Loi-cadre 09-21 : de l'AMO de base à l'Aide Sociale Directe (ASD).
        \item \textbf{B.} Le rôle pivot des outils de ciblage : Registre National de la Population (RNP) et Registre Social Unifié (RSU).
    \end{itemize}
    \item \textbf{II. Le Défi Majeur de la Soutenabilité Budgétaire et de l'Efficience Opérationnelle}
    \begin{itemize}
        \item \textbf{A.} Les mécanismes de financement : articulation entre cotisations contributives et solidarité budgétaire de l'État.
        \item \textbf{B.} Les réformes indispensables d'accompagnement : refonte du système national de santé et réforme de la Caisse de Compensation.
    \end{itemize}
\end{itemize}

\subsection*{3. Données Clés}
\begin{itemize}
    \item Coût global annuel du chantier : environ 51 milliards de dirhams par an, dont près de 25 milliards financés directement par le budget de l'État.
    \item Calendrier : AMO pour tous fin 2022 (transition du RAMED vers AMO Tadamon), lancement des aides directes fin 2023 (allocations de 500 à 1000 DH/mois par famille ciblée), extension prévue pour les régimes de retraite et l'indemnité pour perte d'emploi.
\end{itemize}
